\section{The strongest objections}\label{sec:objections}
\begin{enumerate}
\item \textbf{CBT has a large evidence base.} Yes (Section~\ref{sec:evidence}). The argument concerns mechanism and moderators, not average efficacy.
\item \textbf{CBT for adult ADHD works.} It beats matched controls in medicated adults with a skills-heavy protocol; the evidence is graded low to moderate (Section~\ref{sec:adhdcbt}).
\item \textbf{Patients practise when calm.} Accepted (Section~\ref{sec:ingredients}); it relocates the active ingredient to practice.
\item \textbf{Exposure is behavioural.} Likewise; it strengthens the argument against restructuring specifically, though exposure may act via changed expectancies \citep{hofmann2008extinction}.
\item \textbf{Mindfulness is not better.} True. The argument needs only that restructuring is not uniquely necessary.
\item \textbf{Cognitive therapy sometimes wins.} Yes, in social phobia \citep{ougrin2011exposure} and in one PTSD trial \citep{bryant2008cr}.
\item \textbf{The neuroscience is about acute stress.} Correct, and the human reappraisal-under-stress findings are mixed (Section~\ref{sec:stress}).
\item \textbf{Executive function does not predict CBT outcome.} Largely true for overall outcome (Section~\ref{sec:moderators}); the argument survives only in the narrower, skill-level form.
\item \textbf{Therapists scaffold the executive load.} Plausible and untested; a therapist who generates the evidence list lowers demand and could erase a moderator effect.
\end{enumerate}
\textbf{Limitations.} Narrative review by a non-clinician; single-reviewer screening; several key studies are in older adults or children; evidence from depression stands in for anxiety on some questions (component trials, bias adjustment, decline over time).
